\subsection{Different fixed covariances}\label{sec:different}

This subsection removes the common-covariance assumption. The linear
part remains common, and each component's covariance remains
independent of the state. The branches are then distinct Gaussian
integral operators rather than weighted composition operators composed
with one smoothing operator, so the Gram matrix is no longer fixed up to
unitary conjugation; a compactness argument replaces it.

\begin{lemma}[Uniform independence of Gaussian profiles]\label{lem:uniformgram}
Fix positive definite matrices $S_j$ and vectors $b_j$,
$1\le j\le N$, with pairwise distinct pairs $(S_j,b_j)$.
Let $g_j$ be the $L^2(\R^m)$-normalised Gaussian proportional
to $e^{-\xi^\top S_j\xi/8}$. There exists $\gamma>0$, depending
only on these fixed parameters, such that the Gram matrix of
\[
 u_j(\xi)=g_j(\xi-c_j)e^{i\xi\cdot b_j}
\]
has smallest eigenvalue at least $\gamma$, for all choices
of real centres $c_1,\ldots,c_N$.
\end{lemma}

\begin{proof}
First, any fixed collection of these functions is linearly
independent. Expanding their exponents gives quadratic
coefficients $S_j/8$ and complex linear coefficients
$S_jc_j/4+ib_j$. Choose a real direction $v$ outside the
finitely many nonzero quadratic zero sets
$v^\top(S_i-S_j)v=0$, and, when $S_i=S_j$, outside
$(b_i-b_j)\cdot v=0$. On the complex line $\xi=tv$,
the exponential functions have distinct quadratic coefficients,
or distinct complex linear coefficients when the quadratics agree.
Choose a ray $t=s e^{i\theta}$ with $\cos(2\theta)>0$ and
with distinct real parts of the linear coefficients within
each equal-quadratic group. Only finitely many angles are
excluded. In a purported linear relation, the term with
smallest quadratic coefficient and largest corresponding
linear real part dominates as $s\to\infty$. Its coefficient
is zero. Repeating removes every term. A relation that holds
only almost everywhere on $\R^m$ first extends everywhere
by continuity, and then to the entire functions by uniqueness.

For uniformity, suppose a sequence of centre configurations
made the smallest Gram eigenvalue tend to zero. Pass to a
subsequence so that every pair of centre differences is
bounded, or its norm tends to infinity. Bounded difference
defines finitely many clusters. Within each cluster subtract
one anchor centre and pass to a further subsequence so that
all relative centres converge. Translating the cluster
preserves its Gram matrix up to diagonal unitary phases.
Remove those phases.

Across different clusters the inner products tend to zero:
their absolute values are bounded by the overlap of two
fixed-shape positive Gaussians, at most
$C e^{-c|c_i-c_j|^2}$. Within each cluster the Gram matrix
converges to a fixed Gaussian Gram matrix, positive definite
by the preceding independence argument. The full limiting
matrix is block diagonal with positive definite blocks.
Its smallest eigenvalue is positive, a contradiction.
\end{proof}

\begin{theorem}[Finite mixtures with different covariances]\label{thm:unequal}
Let the pairs $(Q_j,b_j)$ be pairwise distinct, with
$Q_j\succeq0$ and $b_j\in\R^m$. For Borel probabilities
$p_j$ summing to one, define
\[
 Tf(x)=\sum_{j=1}^N p_j(x)\E f(Mx+b_j+G_j),
 \qquad G_j\sim N(0,Q_j).
\]
This operator maps $\HH_P(\R^m)$ boundedly into $\HH_R(\R^d)$
if and only if $M$ has full column rank and every $p_j$
extends entire with
\begin{equation}\label{eq:unequalcriterion}
 \sup_{z=x+iy}|p_j(z)|e^{-2y^\top D_jy}<\infty,\qquad
 D_j=R^{-1}-M^\top(P+4Q_j)^{-1}M.
\end{equation}
For every nonzero weight, $D_j\succeq0$, and that weight is
constant along $\ker D_j$. A branch with a negative direction
in $D_j$ must have identically zero weight.
\end{theorem}

\begin{proof}
The positive-probe argument used for Theorem \ref{thm:mixture}
again makes $M$ injective. Boundedness gives an entire
$L^2(\mu_P)$-valued source representer $V_z$ with
$\norm{V_z}\le\norm T e^{2y^\top R^{-1}y}$.
The pure component features are
\[
 h_j(z)(\xi)=
 e^{-\xi^\top Q_j\xi/2+i\xi\cdot(Mz+b_j)}.
\]
They are entire Hilbert-space-valued functions, and at every
complex $z$ they are linearly independent. Their common
nonvanishing exponential factor cancels, after which the
independence proof in Lemma \ref{lem:uniformgram} applies.

On real states $V_x=\sum_jp_j(x)h_j(x)$. The continuous
multilinear exterior product
$V_z\wedge h_1(z)\wedge\cdots\wedge h_N(z)$ is holomorphic
and vanishes on $\R^d$, hence everywhere. Equivalently,
apply arbitrary $N+1$ bounded linear functionals and continue
the resulting scalar determinants. Independence implies
$V_z\in\operatorname{span}\{h_j(z)\}$ for every complex $z$.
Near a chosen $z_0$, select $N$ bounded functionals whose
value matrix on the $h_j(z_0)$ is invertible. The same
matrix remains invertible nearby; solving for the
coefficients makes the $p_j$ holomorphic there. Uniqueness
patches these coefficients into entire functions agreeing
with the original Borel probabilities.

Put $S_j=P+4Q_j$. In ordinary $L^2(\dd\xi)$ coordinates,
remove the common unitary factor $e^{i\xi\cdot Mx}$.
Completing the square shows that each normalised feature
is a function in Lemma \ref{lem:uniformgram}, with
centre $c_j(y)=-4S_j^{-1}My$ and fixed modulation $b_j$.
Its unnormalised squared norm is
\[
 \norm{h_j(z)}^2=
 \sqrt{\frac{\det P}{\det S_j}}\,
       e^{4y^\top M^\top S_j^{-1}My}.
\]
The uniform Gram lower bound therefore gives
\[
 \gamma\sum_j|p_j(z)|^2
       \sqrt{\frac{\det P}{\det S_j}}\,
       e^{4y^\top M^\top S_j^{-1}My}
 \le\norm{V_z}^2
 \le\norm T^2e^{4y^\top R^{-1}y}.
\]
This proves \eqref{eq:unequalcriterion}. Conversely, set
$E_j=(M^\top S_j^{-1}M)^{-1}$. Each pure component maps
$\HH_P$ boundedly into $\HH_{E_j}$, and
Lemma \ref{lem:multiplier} makes multiplication by $p_j$
a bounded map to $\HH_R$. Their sum is bounded.

Finally, if $v^\top D_jv<0$, the bound on the entire
function $t\mapsto p_j(z_0+tv)$ is a negative quadratic
in $\im t$, plus a linear term, and is independent of
$\re t$. It is bounded on $\C$ and tends to zero as
$|\im t|\to\infty$. Liouville's theorem makes it identically
zero. If $D_j\succeq0$ and $v\in\ker D_j$, the same
argument without decay makes the function constant.
\end{proof}

Linear independence of finitely many distinct Gaussian profiles is
classical, in the same circle of ideas as identifiability of finite
Gaussian mixtures. The new statement is the classification
\eqref{eq:unequalcriterion}. The compactification proof supplies
uniform control over all states. Its constant $\gamma$ is
qualitative. The common-covariance formula \eqref{eq:unitarygram}
provides a simpler explicit constant when that restriction is
available.
